\section{Introduction}
\label{sec:intro}

The cost of searching a binary search tree is the length of the path the search
walks, so the height of the tree determines the cost of every operation the tree
supports. A tree of $n$ nodes can have any height between $\lfloor \log_2 n
\rfloor$ and $n-1$, and the search tree property does not prefer one over the
other. Both extremes are legal. Which one a program ends up with depends on the
order the keys happened to arrive in, and that order is usually not under the
programmer's control.

\Cref{fig:degenerate} shows how little it takes to reach the bad extreme. Both
trees hold the keys 1 through 7 and both obey the search tree property. The tree
on the left was built from an order that happened to interleave small and large
keys, and it answers a query in at most three comparisons. The tree on the right
was built from the keys in increasing order. Every key was larger than every key
before it, so every key became a right child, and the result is a list. Sorted
input is not an exotic case. It is what a program sees when it loads records from
an indexed file, replays a log, or inserts timestamps.

\begin{figure}[t]
  \centering
  % Two legal binary search trees over the same seven keys. The arrival order is
% the only difference between them, and it already costs a factor of more than
% two in search depth at this size.
%
% Positions are explicit rather than delegated to TikZ's tree layout: the right
% tree has one child per node, and the automatic layout would stack those
% vertically instead of showing the list running down to the right.
\begin{tikzpicture}[x = 1mm, y = 1mm]

  % ------------------------------------------------ left: a balanced arrival
  \begin{scope}
    \node[tnode] (a4) at (  0,   0) {4};
    \node[tnode] (a2) at (-14, -13) {2};
    \node[tnode] (a6) at ( 14, -13) {6};
    \node[tnode] (a1) at (-21, -26) {1};
    \node[tnode] (a3) at ( -7, -26) {3};
    \node[tnode] (a5) at (  7, -26) {5};
    \node[tnode] (a7) at ( 21, -26) {7};
    \foreach \p/\c in {a4/a2, a4/a6, a2/a1, a2/a3, a6/a5, a6/a7}
      \draw[tedge] (\p) -- (\c);

    \node[tannot] at (0,  10) {inserted 4, 2, 6, 1, 3, 5, 7};
    \node[tannot] at (0, -36) {height 2, at most three comparisons};
  \end{scope}

  % -------------------------------------------- right: the keys arrive sorted
  \begin{scope}[xshift = 72mm]
    \foreach \k [count = \i from 0] in {1, 2, 3, 4, 5, 6, 7}
      \node[tnode] (b\k) at (\i * 6 - 18, -\i * 8) {\k};
    \foreach \p/\c in {b1/b2, b2/b3, b3/b4, b4/b5, b5/b6, b6/b7}
      \draw[tedge] (\p) -- (\c);

    \node[tannot] at (0,  10) {inserted 1, 2, 3, 4, 5, 6, 7};
    \node[tannot] at (0, -60) {height 6, up to seven comparisons};
  \end{scope}

\end{tikzpicture}

  \caption{Two binary search trees over the same seven keys. Only the insertion
  order differs. The tree on the right is a list, and searching it costs a
  comparison per key.}
  \label{fig:degenerate}
\end{figure}

An AVL tree removes the choice. It keeps the ordinary search tree property and
adds one local requirement: at every node, the heights of the two subtrees differ
by at most one. The requirement is checked and repaired as part of every
insertion and every deletion, so no sequence of updates can produce the tree on
the right of \Cref{fig:degenerate}. Adelson-Velsky and Landis introduced the
structure in 1962~\cite{avl1962}, and it was the first search tree to maintain
its own balance.

The requirement is stated one node at a time, and it is worth asking what a
condition that local can buy globally. The answer is the subject of
\Cref{sec:height}. A tree in which every node is nearly balanced cannot be tall,
because the smallest tree that reaches a given height grows in proportion to
$\varphi^h$, where $\varphi = (1+\sqrt{5})/2$ is the golden ratio. Height is
therefore logarithmic in the number of keys, with a constant of about $1.44$
against the perfect tree.

\subsection*{Contributions}

This report extends a presentation the authors gave on the same subject. It
contributes the following.

\begin{enumerate}
  \item A derivation of the height bound of \Cref{thm:heightbound} that carries
    each step, from the recurrence for the minimum-size tree through the
    characteristic equation to the constant $1.4404$, together with a remark on
    the two height conventions that circulate in the literature and the index
    shift between them (\Cref{rem:conventions}).
  \item Proofs that rebalancing preserves the in-order key sequence
    (\Cref{lem:inorder}) and that it costs constant time
    (\Cref{lem:rotcost}), with the four repair cases and the algorithms that
    apply them.
  \item An analysis of why one insertion needs at most one rebalancing step
    (\Cref{thm:insertone}) while one deletion may need a rebalancing step at
    every level along the path to the root (\Cref{thm:deletebound}), and a
    worked minimum-size tree in which three of them fire.
  \item Measurements on trees of up to one million keys, reported in
    \Cref{sec:eval}, covering height against the analytic bounds, search cost
    against an unbalanced tree, and rebalancing work per operation.
  \item A comparison with other balanced search trees
    (\Cref{sec:comparison}) and an account of where the structure is used in
    practice (\Cref{sec:applications}).
\end{enumerate}

\Cref{sec:prelim} fixes notation and states the invariant.
\Cref{sec:height} proves the height bound. \Cref{sec:rotations} introduces
rotation as the repair operation, and \Cref{sec:insertion} and
\Cref{sec:deletion} apply it to the two update operations.
\Cref{sec:eval} reports the measurements, \Cref{sec:comparison} sets the
structure against its alternatives, \Cref{sec:applications} describes where it
is used, and \Cref{sec:history} returns to the 1962 paper.
